\chapter{学情建模}

\section{自述解析}

学习者输入一段自我描述，系统从中抽取学历、阶段、专业方向与实操学时。
配置了模型接口时走模型抽取再用规则校验，未配置时纯规则运行，离线一样可用。

需要强调的是，\emph{抽取结果只用来给初始掌握概率定先验，不直接决定结论}。
先验会随着作答证据的累积被冲淡，因此即使抽错也不会一错到底。
这是有意的设计：自述是主观的，学习者容易高估或低估自己，不能让它拍板。

\begin{keynote}{踩坑：「二年级」被算成两年工龄}
自述“高职电气自动化\emph{二年级}”当中的“二年”曾被正则表达式当作时长，
折合三千二百小时实操，先验直接顶到取值上界，测评开局就问偏了难度。
学制与工龄必须分开，现在的做法是在时长模式后加负向前瞻，
排除“级”“段”“检”等后缀。
\end{keynote}

\section{贝叶斯知识追踪}

掌握度估计采用贝叶斯知识追踪\cite{corbett1994}，而非简单的答对率。
答对率有两处硬伤。其一是\emph{蒙对算掌握}：四选一题目的蒙对概率为四分之一，
答对率把这部分当成了真会。其二是\emph{失误算不会}：真正掌握的人手滑点错，
答对率会直接把他判成盲区。

贝叶斯知识追踪用四个参数把这两件事显式建模，各参数含义见表 \ref{tab:bkt}。

\begin{table}[htbp]\centering\small
\caption{贝叶斯知识追踪的四个参数}\label{tab:bkt}
\begin{tabular}{@{}p{1.8cm}p{3.4cm}p{6.2cm}@{}}
\toprule
参数 & 含义 & 说明 \\
\midrule
$p_{L0}$ & 初始掌握概率 & 由学历与实操学时给出先验 \\
$p_T$ & 学习跃迁概率 & 前测阶段无教学干预，影响很小 \\
$p_S$ & 失误率 & 已掌握却答错 \\
$p_G$ & 蒙对率 & 四选一题目的理论下界即为 0.25 \\
\bottomrule
\end{tabular}
\end{table}

设 $P(L)$ 为当前掌握概率，一次作答后的后验与学习跃迁按下式更新：

\begin{align}
P(L \mid \text{答对}) &= \frac{P(L)(1-p_S)}{P(L)(1-p_S) + (1-P(L))\,p_G} \\[2pt]
P(L \mid \text{答错}) &= \frac{P(L)\,p_S}{P(L)\,p_S + (1-P(L))(1-p_G)} \\[2pt]
P(L') &= P(L \mid \text{obs}) + \bigl(1 - P(L \mid \text{obs})\bigr)p_T
\end{align}

\begin{keynote}{必须卡住的一条：$p_S + p_G < 1$}
这两个参数之和大于等于一时模型退化，答对与答错所给出的信息方向会\emph{翻转}，
出现“答错反而涨掌握度”的荒唐结果，文献中称为模型退化。
代码中的参数校验函数会直接抛出异常拦住此类配置。
\emph{用真实作答数据重新拟合参数时，这一条必须先卡。}
\end{keynote}

\section{证据强度：点估计不足以下结论}

前一节给出的是掌握概率的\emph{点估计}。只看点估计会得出荒唐的结论，
表 \ref{tab:luck} 列出了几种典型作答。

\begin{table}[htbp]\centering\small
\caption{点估计与蒙对概率的对照（四选一题）}\label{tab:luck}
\begin{tabular}{@{}p{1.6cm}p{2.2cm}p{2.6cm}p{5cm}@{}}
\toprule
作答 & 点估计 & 纯蒙达到的概率 & 80\% 区间下界 \\
\midrule
2/2 & 0.896 & 6.2\% & 0.33 \\
3/3 & 0.973 & 1.6\% & 0.48 \\
3/4 & 0.856 & 5.1\% & 0.26 \\
4/4 & 0.994 & 0.4\% & 0.59 \\
\bottomrule
\end{tabular}
\end{table}

原实现把 2/2 判为“掌握牢固”。但纯靠蒙达到 2/2 的概率有 6.2\%，
二十个人里就有一个能蒙出来；3/4 判“掌握牢固”，蒙对概率 5.1\%，同样站不住。
\emph{少量作答下的点估计本身就是高方差的，拿它当结论，
等于把“目前的最佳猜测”说成“已经确认”。}

因此每个知识点除点估计外，必须同时给出两个数。
\emph{蒙对概率}即完全不会的人靠蒙达到该成绩或更好的概率，
按二项分布右尾计算，这是最直观的一个数，
直接回答“我瞎蒙也能考成这样吗”。
\emph{可信区间}即掌握概率的区间估计，其下界才是能对外声称的下限。

区间的算法是：作答正确率 $\theta$ 的后验取 $\mathrm{Beta}(1+k,\,1+n-k)$，
再按观测模型反解回掌握概率：

\begin{equation}
\theta = p_L(1-p_S) + (1-p_L)p_G
\quad\Longrightarrow\quad
p_L = \frac{\theta - p_G}{1 - p_S - p_G}
\end{equation}

这一步反解正是\emph{扣掉蒙对成分}。分母必须为正，
也就是 $p_S + p_G < 1$，否则反解方向会翻转。

区间水平取 $0.80$ 而非常用的 $0.95$：本项目的作答数极少，
$0.95$ 区间在 $n=2$ 时几乎覆盖全域，虽然正确但没有任何决策价值。

\subsection{判定门槛}

声称“已确认掌握”要同时满足三条：点估计过线、\emph{区间下界也过线}、
且蒙对概率不高于 5\%。三条里最容易被忽略的是第二条——
点估计 $0.896$ 而下界只有 $0.33$ 时，说“已掌握”是没有依据的。
不满足则一律降级为“疑似掌握”并标注证据不足。

判盲区不必卡蒙对概率——蒙对只会把成绩往上抬，
成绩低说明连蒙都没蒙上，结论方向是安全的；但一题作答仍不足以定论。

\begin{keynote}{一个必须正视的结论：要确认掌握需要八题全对}
按上述门槛推算，$80\%$ 区间下界要过 $0.80$ 这条线，需要连续八题全对。

这意味着\emph{现在每知识点三道题的题库，结构上就不可能“确认”任何一个知识点}。
系统当前对所有知识点给出的都是“疑似掌握”或“证据不足”。
这不是模型效果问题，是题量问题。
把题库扩到每知识点五到六道之所以列为最高优先级之一，原因即在于此。

在此之前，界面与报告中一律使用“疑似掌握”，不得写“已掌握”。
\end{keynote}

\section{自适应选题}

系统不会把全部题目一次性发给学习者，而是每次只发一道，
根据已有作答实时挑选下一道。设当前掌握概率为 $p_L$，则预测答对概率为

\begin{equation}
P(\text{答对}) = p_L(1-p_S) + (1-p_L)\,p_G
\end{equation}

该值越接近 $0.5$，题目的信息量越大——有把握答对或答错的题目，
问了也不会带来新信息。因此选题时取 $|P - 0.5|$ 最小者，
再叠加一个“该知识点问得越少越优先”的权重。

\section{动态追问}

动态追问直接对应榜题创新性条款中“大模型驱动下的动态追问”这一要求。
当一次作答与已有证据明显冲突时，系统就地追加同一知识点的一道题目，
两类触发情形见表 \ref{tab:probe}。

\begin{table}[htbp]\centering\small
\caption{动态追问的触发条件}\label{tab:probe}
\begin{tabular}{@{}p{4.4cm}p{4cm}p{3cm}@{}}
\toprule
情形 & 歧义 & 动作 \\
\midrule
估计已掌握（不低于 0.65）却答错 & 是失误还是真不会 & 追问一次 \\
估计为盲区（不高于 0.35）却答对 & 是真会还是蒙对 & 追问一次 \\
\bottomrule
\end{tabular}
\end{table}

贝叶斯知识追踪的失误率与蒙对率恰好把这两件事显式建模了，因而歧义可以量化。
\emph{追问不是随便加一道题，而是为消解特定歧义而加。}

\begin{keynote}{踩坑：追问在第一题就触发}
最初没有限制触发前提，结果先验偏低时（零基础学员约 0.20），
\emph{任何一道答对都满足“估计为盲区却答对”}，第一题就触发追问。
但那时手上一条证据都没有，所谓冲突是与一个通用先验冲突，消解不了任何歧义。
实测十六道题中有七道是追问，知识点覆盖率被吃掉近一半。
修法是加一道门槛：\emph{该知识点必须已有作答证据}才允许追问。
\end{keynote}

\section{能力图谱}

逐条列出十五个知识点的掌握概率，信息是全的，但没人读得进去——
一屏数字，看不出这个人强在哪、弱在哪。
因此把知识点聚合成六个能力维度，以雷达图呈现。

雷达图有个众所周知的毛病：它很好看，因而特别容易骗人。
本系统用三条设计压住这一点。

\emph{轴不能太碎。} 不直接用知识点标签做轴——
现有十三个标签里有八个只覆盖一个知识点，
画出来大部分轴由单点决定，一道题的对错就能让整条轴塌掉，
看着像能力图谱，其实是噪声图。归并到六维、每维至少三个知识点才有平均的意义。

\emph{必须画两层。} 外层是点估计，内层是区间下界，
两层之间的间隙就是“还不确定的部分”。
只画点估计会让四道题答对三道看起来和真正掌握一模一样，
而这两件事的证据强度差着数量级。

\emph{未测的维度不能画成零。} 零和“没测过”在雷达图上长得一样，但含义相反：
前者是“确定不会”，后者是“不知道”。
未测部分单独用虚线标出，并在图例中写明。

六个维度分别为基础认知、现场操作、安全规范、精度标定、编程集成与运维调试，
合计覆盖全部十五个知识点。维度得分取该维度下各知识点的平均，
\emph{未测的知识点不计入平均}，而是单独记进覆盖率——
把未测当零会让没考的维度显示成“完全不会”，这是雷达图最常见的误导方式。

\section{前置推断}

十五个知识点若要各自测到可信度 $0.75$，需要每点三题、合计四十五题，
没有学习者愿意作答。因此系统增加了一条规则：
\emph{前置知识点已确认为盲区的，不再实测，直接推断为盲区}。
连坐标系都分不清的学习者不可能掌握工具坐标系四点标定，
花一道题去确认这件事是浪费。

推断得出的盲区在界面上用虚线框标出，\emph{不冒充实测结果}，
报告中亦应分开统计。

\bibliographystyle{gbt7714-2005}
\bibliography{refs}
